% Cardinal Grid Notes: IEEE conference-style template (IEEEtran, two columns).
% Build:  pdflatex paper && bibtex paper && pdflatex paper && pdflatex paper
% Rules:  figures 3.5in (one column) or 7.16in (two columns); no titles inside figures, use captions;
%         impersonal voice; every number produced by the analysis script; every claim with a citation.
\documentclass[conference]{IEEEtran}
\IEEEoverridecommandlockouts
\usepackage[utf8]{inputenc}
\usepackage[T1]{fontenc}
\usepackage{cite}
\usepackage{amsmath,amssymb}
\usepackage{graphicx}
\usepackage{booktabs}
\usepackage{url}
\usepackage[hidelinks]{hyperref}
\usepackage{microtype}
\def\BibTeX{{\rm B\kern-.05em{\sc i\kern-.025em b}\kern-.08em T\kern-.1667em\lower.7ex\hbox{E}\kern-.125emX}}

\begin{document}

\title{Title of the Note\\
\large Cardinal Grid Notes, No.\ N}

\author{\IEEEauthorblockN{Ricardo W. C. Guerra Filho}
\IEEEauthorblockA{Cardinal Grid (independent open initiative)\\
ORCID 0000-0001-6699-9951 \quad contact@cardinalgrid.com}}

\maketitle

\begin{abstract}
Two hundred words at most. Problem, data, method, main quantitative results, limitation, and where the code is.
\end{abstract}

\begin{IEEEkeywords}
load forecasting, balancing authority, EIA-930, forecast evaluation, data quality
\end{IEEEkeywords}

\section{Introduction}
Context and the specific question, with citations to the public record.

\section{Data}
Source, definitions in the data owner's own terms, coverage.

\section{Methods}
Metrics with equations, data-quality rules, eligibility, event windows.

\section{Results}
Figures and tables only; interpretation goes to the Discussion.

\begin{figure}[t]
\centering
\includegraphics[width=\columnwidth]{figures/fig1.pdf}
\caption{Caption written as a full sentence, including the data source and the period.}
\label{fig:one}
\end{figure}

\section{Discussion}
What the results mean, what they do not show, what comes next.

\section{Conclusion}
Three sentences.

\section*{Data and Code Availability}
Public data sources and the archived code with its DOI.

\bibliographystyle{IEEEtran}
\bibliography{references}

\end{document}
